% Fichier généré automatiquement par tools/generate_corrections.py.
% Source : PERF/PERF-05-Precistion-TVF/509_Divers/509_Divers.tex
% Statut : complet (5/5 question(s) renseignée(s), 0 reprise(s)).
\begin{corrigebox}[Corrigé extrait de la banque]
\CorrigeQuestion{1}
On a : 
$\varepsilon(p) = E(p)-C\dfrac{1}{p\left( A+\tau_1 p \right)} \left( \varepsilon(p)\dfrac{K}{p} + P(p)\right)$

$\Leftrightarrow \varepsilon(p) \left(1 +C\dfrac{K}{p^2\left( A+\tau_1 p \right)}\right) = E(p)-C\dfrac{1}{p\left( A+\tau_1 p \right)}  P(p)$

$\Leftrightarrow \varepsilon(p) = E(p)\dfrac{1}{1 +C\dfrac{K}{p^2\left( A+\tau_1 p \right)}}-C\dfrac{1}{p\left( A+\tau_1 p \right)} \dfrac{1}{1 +C\dfrac{K}{p^2\left( A+\tau_1 p \right)}} P(p)$

$\Leftrightarrow \varepsilon(p) = E(p)\dfrac{1}{1 +C\dfrac{K}{p^2\left( A+\tau_1 p \right)}}-\dfrac{C}{p\left( A+\tau_1 p \right) +C\dfrac{K}{p}} P(p)$
\CorrigeQuestion{2}
On a $\lim\limits_{t\to +\infty} \varepsilon(t) = \lim_{p\to 0} p\varepsilon(p)$.

Dans ces conditons, 
$\varepsilon(p) = E_0\dfrac{1}{p +C\dfrac{K}{p\left( A+\tau_1 p \right)}}-\dfrac{C}{p^2\left( A+\tau_1 p \right) +CK} P_0$.

Au final, $\varepsilon = E_0\dfrac{p}{p +C\dfrac{K}{p\left( A+\tau_1 p \right)}}-\dfrac{Cp}{p^2\left( A+\tau_1 p \right) +CK} P_0$ $=0-0 =0$
\CorrigeQuestion{3}
Dans ces conditons, 
$\varepsilon(p) = E_0\dfrac{1}{p +C\dfrac{K}{p\left( A+\tau_1 p \right)}}-\dfrac{C}{p^3\left( A+\tau_1 p \right) +CKp} P_0$

et  $\varepsilon = p^2E_0\dfrac{1}{p^2 +C\dfrac{K}{\left( A+\tau_1 p \right)}}-\dfrac{C}{p^2\left( A+\tau_1 p \right) +CK} P_0 = 0 - \dfrac{P_0}{K}=- \dfrac{P_0}{K}$
\CorrigeQuestion{4}
Dans ces conditons, 
$\varepsilon(p) = E_0\dfrac{1}{p^2 +C\dfrac{K}{\left( A+\tau_1 p \right)}}-\dfrac{C}{p^2\left( A+\tau_1 p \right) +CK} P_0$.

On a alors $\varepsilon = pE_0\dfrac{1}{p^2 +C\dfrac{K}{\left( A+\tau_1 p \right)}}-\dfrac{C}{p^2\left( A+\tau_1 p \right) +CK} P_0p = 0$
\CorrigeQuestion{5}
Dans ces conditons, 
$\varepsilon(p) = E_0\dfrac{1}{p^2 +C\dfrac{K}{\left( A+\tau_1 p \right)}}-\dfrac{C}{p^3\left( A+\tau_1 p \right) +CKp} P_0$. 
On a alors
$\varepsilon = E_0 p \dfrac{1}{p^2 +C\dfrac{K}{\left( A+\tau_1 p \right)}}-\dfrac{C}{p^2\left( A+\tau_1 p \right) +CK} p P_0 = -\dfrac{P_0}{K} $.
\end{corrigebox}
